% Part: first-order-logic
% Chapter: introduction
% Section: models-theories

\documentclass[../../../include/open-logic-section]{subfiles}

\begin{document}

\olfileid{fol}{int}{mod}

\olsection{నమూనాలు మరియు సిద్ధాంతాలు}

మొదటిస్థాయి విధేయ తర్కపు వాక్యనిర్మాణం, అర్థవిచారం నిర్వచించిన తరువాత,
\tetoken{నిర్మాణాలు}{!!{structure}s} ధర్మాలను, అర్థపర భావనలను పరిశీలించడం మొదలుపెట్టవచ్చు.
\tetoken{వ్యుత్పత్తి}{!!{derivation}} వ్యవస్థలను కూడా నిర్వచించి పరిశీలించవచ్చు. ఒక
\tetoken{వాక్యాలు}{!!{sentence}s} సమితి గురించి, ఆ సమితిలోని అన్ని \tetoken{వాక్యాలను}{!!{sentence}s} ఏ
\tetoken{నిర్మాణాలు}{!!{structure}s} సత్యం చేస్తాయి అని అడగవచ్చు. \tetoken{వాక్యాలు}{!!{sentence}s}
సమితి~$\Gamma$ ఇచ్చినప్పుడు, వాటిని సంతృప్తిపరిచే
\tetoken{ఒక నిర్మాణం}{!!a{structure}}~$\Struct{M}$ను \emph{$\Gamma$ యొక్క నమూనా} అంటాం.
$\Gamma$తో మొదలుపెట్టి దాని నమూనాలను కనుగొనడానికి ప్రయత్నించవచ్చు:
అవి ఎలా ఉంటాయి? అవి ఎంత పెద్దవిగా లేదా చిన్నవిగా ఉండాలి? లేదా ఒకే
\tetoken{ఒక నిర్మాణంతో}{!!{structure}} గానీ, \tetoken{నిర్మాణాలు}{!!{structure}s} సమూహంతో గానీ మొదలుపెట్టి,
వాటిలో ఏ \tetoken{వాక్యాలు}{!!{sentence}s} సత్యమని అడగవచ్చు. ఈ \tetoken{నిర్మాణాలను}{!!{structure}s}
\emph{లాక్షణీకరించే}, అంటే వాటిలో సత్యమయ్యే వాక్యాలు సరిగ్గా అవే అని
నిర్ణయించే, \tetoken{వాక్యాలు}{!!{sentence}s} ఉన్నాయా? ఇలాంటి ప్రశ్నలు \emph{నమూనా
సిద్ధాంతం} పరిధిలోకి వస్తాయి. ఒక \tetoken{నిర్మాణాలు}{!!{structure}s} సమూహాన్ని ఒక
\tetoken{వాక్యాలు}{!!{sentence}s} సమితితో, అంటే ఒక సిద్ధాంతపు స్వీకృతాలతో వర్ణించే
\emph{స్వీకృతాధారిత పద్ధతి}కి కూడా ఇవే ఆధారం. స్వీకృతాల నుంచి
మొదటిస్థాయి విధేయ తర్కంలో అనుగమించే \tetoken{వాక్యాలు}{!!{sentence}s} సరిగ్గా ఆ స్వీకృతాల అన్ని
నమూనాల్లోనూ సత్యమవుతాయనే పరిశీలన దీన్ని సాధ్యం చేస్తుంది.

చాలా సరళమైన ఉదాహరణగా పూర్వక్రమాలను పరిశీలించండి. ఏదో సమితి~$A$పై
ఉన్న సంబంధం~$R$ స్వావర్తనమూ సంక్రామకమూ అయితే అది పూర్వక్రమం. ఒక
ద్విస్థాన సంబంధం~$R \subseteq A \times A$తో కూడిన సమితి~$A$,
ఒకే ద్విస్థాన సంబంధ సంకేతం~$\Obj P$ గల మొదటిస్థాయి భాషకు
\tetoken{ఒక నిర్మాణం}{!!a{structure}} ఇవ్వడానికి మనకు కావాల్సినదే: $\Domain{M} = A$,
$\Assign{\Obj P}{M} = R$గా పెడతాం. $R$ పూర్వక్రమం కాబట్టి అది
స్వావర్తనమూ సంక్రామకమూ. ఈ విషయాలను చెప్పే మొదటిస్థాయి విధేయ తర్కపు
\tetoken{వాక్యాలు}{!!{sentence}s} సమితి~$\Gamma$ను కనుగొనవచ్చు:
\begin{align*}
  & \lforall[\Obj v_0][\Atom{\Obj P}{\Obj v_0,\Obj v_0}]\\
  & \lforall[\Obj v_0][\lforall[\Obj v_1][\lforall[\Obj v_2][((\Atom{\Obj P}{\Obj v_0,\Obj v_1} \land \Atom{\Obj P}{\Obj v_1,\Obj v_2}) \lif \Atom{\Obj P}{\Obj v_0,\Obj v_2})]]]
\end{align*}
ఈ \tetoken{వాక్యాలు}{!!{sentence}s}, ``ప్రతి~$x$కు $Rxx$'' ($R$ స్వావర్తనం),
``$Rxy$, $Ryz$ అయినప్పుడల్లా $Rxz$ కూడా'' ($R$ సంక్రామకం) అనే
వాక్యాల సంకేతీకరణలే. \tetoken{ఒక నిర్మాణం}{!!a{structure}}~$\Struct{M}$ ఈ రెండు
\tetoken{వాక్యాలు}{!!{sentence}s} సమితి~$\Gamma$కు నమూనా కావడం, $R$ (అంటే
$\Assign{\Obj P}{M}$) అనేది $A$పై (అంటే $\Domain{M}$పై) ఒక
పూర్వక్రమమైనప్పుడూ, అప్పుడు మాత్రమే అని చూస్తాం. మరోలా చెప్పాలంటే,
$\Gamma$ యొక్క నమూనాలు సరిగ్గా పూర్వక్రమాలే. కేవలం~$\Obj P$ను
\tetoken{ఒక విధేయ సంకేతంగా}{!!{predicate}} కలిగిన మొదటిస్థాయి భాషలో వ్యక్తీకరించగల అన్ని
పూర్వక్రమాల ఏ ధర్మమైనా (పైనున్న స్వావర్తనత్వం, సంక్రామకత్వంలా),
$\Gamma$లోని రెండు \tetoken{వాక్యాలు}{!!{sentence}s} నుంచి అనుగమిస్తుంది; దాని విపర్యయం
కూడా నిలుస్తుంది. కాబట్టి $\Gamma$ నమూనాల గురించి మనం నిరూపించగల
ప్రతి విషయాన్నీ అన్ని పూర్వక్రమాల గురించి నిరూపించినట్లే.

ఏదైనా నిర్దిష్ట సిద్ధాంతం, దాని నమూనాల వర్గం ($\Gamma$, అన్ని
పూర్వక్రమాలు వంటివి) కోసం, సంబంధిత మొదటిస్థాయి భాషలో దేన్ని
వ్యక్తీకరించవచ్చు, దేన్ని వ్యక్తీకరించలేం అనే ఆసక్తికరమైన ప్రశ్నలు
ఉంటాయి. అన్ని భాషలు, నమూనాల వర్గాలకు ఆసక్తికరమైన కొన్ని
\tetoken{నిర్మాణాలు}{!!{structure}s} ధర్మాలు ఉన్నాయి; అవి \tetoken{వ్యక్తి క్షేత్రం}{!!{domain}} పరిమాణానికి సంబంధించినవి.
ఉదాహరణకు, ఏ~$n \in \PosInt$కైనా \tetoken{వ్యక్తి క్షేత్రంలో}{!!{domain}} సరిగ్గా
$n$~\tetoken{మూలకాలు}{!!{element}s} ఉన్నాయని ఎల్లప్పుడూ వ్యక్తీకరించవచ్చు. అనంత
సంఖ్యలో \tetoken{వాక్యాలు}{!!{sentence}s} గల సమితిని ఉపయోగించి, \tetoken{వ్యక్తి క్షేత్రం}{!!{domain}} అనంతమని కూడా
వ్యక్తీకరించవచ్చు. కానీ ప్రదేశం పరిమితమని గానీ, అది
\tetoken{లెక్కించలేని}{!!{nonenumerable}} అని గానీ వ్యక్తీకరించలేం. మొదటిస్థాయి భాషల ఈ
పరిమితుల గురించిన ఫలితాలు సంహతత్వ, లొవెన్‌హైమ్--స్కోలెమ్ సిద్ధాంతాల
పరిణామాలు.

\end{document}
